\section*{Bab \ref{ch:b1:mirrors-thin-lenses} --- Cermin dan Lensa Tipis}

\begin{solution}{exo:b1:mirrors-thin-lenses:1}
$\overline{SF} = \qty{-20}{cm}$; $1/\overline{SA'} = 1/\overline{SF} -
1/\overline{SA}$. Pada \qty{60}{cm}: $\overline{SA'} = \qty{-30}{cm}$,
nyata, $\gamma = -0.5$. Pada \qty{30}{cm}: $\overline{SA'} = \qty{-60}{cm}$,
nyata, $\gamma = -2$. Pada \qty{10}{cm}: $1/\overline{SA'} = -1/20 + 1/10
= 1/20$, $\overline{SA'} = \qty{+20}{cm}$ di belakang cerminnya, maya,
$\gamma = +2$ (tegak, diperbesar).
\end{solution}

\begin{solution}{exo:b1:mirrors-thin-lenses:2}
$\overline{OA} = \qty{-20}{cm}$: $1/\overline{OA'} = 1/5 - 1/20 = 3/20$,
$\overline{OA'} = \qty{+6.7}{cm}$, nyata terbalik, $\gamma = -1/3$.
$\overline{OA} = \qty{-3.0}{cm}$: $1/\overline{OA'} = 1/5 - 1/3 = -2/15$,
$\overline{OA'} = \qty{-7.5}{cm}$, maya tegak, $\gamma = +2.5$.
\end{solution}

\begin{solution}{exo:b1:mirrors-thin-lenses:3}
$V = 1/0.25 = \qty{+4.0}{\delta}$; $f' = 1/(-2.0) = \qty{-50}{cm}$;
$V = 3.0 - 1.0 = \qty{+2.0}{\delta}$, $f' = \qty{50}{cm}$.
\end{solution}

\begin{solution}{exo:b1:mirrors-thin-lenses:4}
$\overline{FA} = \qty{-15}{cm}$: $\overline{F'A'} = -f'^2/\overline{FA}
= +100/15 = \qty{6.7}{cm}$ melewati $F'$, yaitu $\overline{OA'} = \qty{16.7}{cm}$;
$\gamma = f'/\overline{FA} = -0.67$. Descartes: $\overline{OA} =
\qty{-25}{cm}$, $1/\overline{OA'} = 1/10 - 1/25 = 3/50$,
$\overline{OA'} = \qty{16.7}{cm}$, $\gamma = 16.7/(-25) = -0.67$.
\end{solution}

\begin{solution}{exo:b1:mirrors-thin-lenses:5}
$\overline{SF} = \qty{+1.0}{m}$, $\overline{SA} = \qty{-20}{m}$:
$1/\overline{SA'} = 1 + 1/20$, $\overline{SA'} = \qty{0.95}{m}$ di
belakang cerminnya (maya, tegak); $\gamma = -0.95/(-20) = +0.048$,
tinggi \omterm{def:b1:mirrors-thin-lenses:image}{bayangannya} \qty{7.1}{cm}. Otak menilai jarak dari ukuran
semunya: mobil yang tampak dua puluh kali lebih kecil pada \qty{0.95}{m}
dibaca sebagai mobil yang jauh --- lebih jauh daripada \qty{20}{m}.
\end{solution}

\begin{solution}{exo:b1:mirrors-thin-lenses:6}
$1/\overline{OA'} = 1/50 - 1/2000 = \qty{0.0195}{mm^{-1}}$,
$\overline{OA'} = \qty{51.3}{mm}$: \omterm{def:b1:mirrors-thin-lenses:lens}{lensanya} bergerak \qty{1.3}{mm} menjauhi
sensornya. $\gamma = 51.3/(-2000) = -0.026$. Untuk $\abs\gamma = 24/1800
= 1/75$: $\overline{OA} = f'(1/\gamma - 1) = 50 \times (-76) =
\qty{-3.8}{m}$.
\end{solution}

\begin{solution}{exo:b1:mirrors-thin-lenses:7}
$\gamma = -1.8/0.036 = -50$; $\overline{OA'} = \qty{5.0}{m}$,
$\overline{OA} = \overline{OA'}/\gamma = \qty{-10}{cm}$;
$1/f' = 1/5.0 + 1/0.10 = \qty{10.2}{m^{-1}}$, $f' = \qty{98}{mm}$.
\omterm{def:b1:mirrors-thin-lenses:image}{Bayangannya} terbalik ($\gamma < 0$), atas--bawah dan kiri--kanan: slidenya
dimasukkan setelah diputar $180^\circ$.
\end{solution}

\begin{solution}{exo:b1:mirrors-thin-lenses:8}
Benda di bidang fokus $\Rightarrow$ sinar dari tiap titiknya keluar dari
\omterm{def:b1:mirrors-thin-lenses:lens}{lensanya} sejajar; \omterm{prop:b1:rays-reflection-refraction:mirror}{cermin datarnya} mengirimkannya kembali sejajar (arah
setangkup); \omterm{def:b1:mirrors-thin-lenses:lens}{lensanya} memfokuskan berkas sejajar di bidang fokusnya, pada
titik yang setangkup terhadap sumbu dari titik sumbernya: \omterm{def:b1:mirrors-thin-lenses:image}{bayangannya} di
bidang bendanya, terbalik, $\gamma = -1$. Memiringkan cerminnya sebesar
$\alpha$ memutar berkas yang kembali sebesar $2\alpha$ dan menggeser
\omterm{def:b1:mirrors-thin-lenses:image}{bayangannya} ke samping sebesar $2\alpha f'$, tetap tajam di bidang yang sama.
\end{solution}

\begin{solution}{exo:b1:mirrors-thin-lenses:9}
$f' = (100^2 - 40^2)/(4 \times 100) = \qty{21}{cm}$. Jarak terpendeknya
$4f' = \qty{84}{cm}$, dengan $\gamma = -1$.
\end{solution}

\begin{solution}{exo:b1:mirrors-thin-lenses:10}
$L_1$ membayangkan tak hingga di $F_1'$, \qty{20}{cm} melewati $L_1$, yaitu \qty{5}{cm}
melewati $L_2$: sebuah benda maya, $\overline{O_2A_1} = \qty{+5}{cm}$.
$1/\overline{O_2A'} = -1/10 + 1/5 = 1/10$: \omterm{def:b1:mirrors-thin-lenses:image}{bayangan} akhirnya \qty{10}{cm}
melewati $L_2$. Tinggi \omterm{def:b1:mirrors-thin-lenses:image}{bayangan} antaranya $h_1 = f_1'\alpha = 20\alpha$ (cm);
$\gamma_2 = 10/5 = 2$, sehingga $h = 40\alpha$: \omterm{def:b1:mirrors-thin-lenses:lens}{jarak fokus} setaranya
\qty{40}{cm}, dalam tabung sepanjang $15 + 10 = \qty{25}{cm}$ --- sebuah telefoto.
\end{solution}

\begin{solution}{exo:b1:mirrors-thin-lenses:11}
$\gamma = -3 = -\overline{SA'}/\overline{SA}$ memberi $\overline{SA'} =
3\overline{SA}$; keduanya nyata, di depan: $\overline{SA} = -x$,
$\overline{SA'} = -3x$, dan jarak layar--lampu $2x = \qty{80}{cm}$:
$x = \qty{40}{cm}$. Cerminnya \qty{40}{cm} dari lampunya, \qty{120}{cm}
dari layarnya; $1/\overline{SF} = -1/40 - 1/120 = -1/30$:
$R = 2 \times 30 = \qty{60}{cm}$.
\end{solution}

\begin{solution}{exo:b1:mirrors-thin-lenses:12}
$\theta = 0.53^\circ = \qty{9.3e-3}{rad}$; garis tengah \omterm{def:b1:mirrors-thin-lenses:image}{bayangannya} $f'\theta =
\qty{4.6}{mm}$. Daya yang terkumpul $10^3 \times \pi(0.05)^2 = \qty{7.9}{W}$;
luas \omterm{def:b1:mirrors-thin-lenses:image}{bayangannya} $\pi(2.3\times10^{-3})^2 = \qty{1.7e-5}{m^2}$; iradiansinya
$\qty{4.7e5}{W/m^2}$, $470$ kali sinar matahari langsung: kertasnya
menghangus lalu terbakar --- kaca pembakar. Secara umum penguatannya
$(D/f'\theta)^2$.
\end{solution}

\begin{solution}{pb:b1:mirrors-thin-lenses:1}
\textbf{1.} $1/\overline{OA'} - 1/\overline{OA} = 1/f'$,
$\gamma = \overline{OA'}/\overline{OA}$; sumbunya diarahkan searah
cahayanya, jaraknya aljabar dari $O$.

\textbf{2.} $p = -\overline{OA}$, $\overline{OA'} = D - p$:
$1/(D - p) + 1/p = 1/f_1'$, yaitu $p^2 - Dp + Df_1' = 0$; akarnya nyata jika dan hanya jika
$D^2 - 4Df_1' \geq 0$, yaitu $D \geq 4f_1'$.

\textbf{3.} Kedua akarnya berjumlah $D$: $p_1 + p_2 = D$, sehingga $p_2 = D - p_1$ ---
jarak benda dan jarak \omterm{def:b1:mirrors-thin-lenses:image}{bayangannya} bertukar, dan kedudukannya setangkup
terhadap titik tengahnya. $d = p_2 - p_1 = \sqrt{D^2 - 4Df_1'}$, jadi $f_1' = (D^2 -
d^2)/4D$.

\textbf{4.} $f_1' = (6400 - 1600)/320 = \qty{15.0}{cm}$.

\textbf{5.} Hanya perpindahan $d$ tabungnya antara kedua setelan tajam
itu dan jarak benda--layar $D$ yang masuk hitungan: keduanya tak
menuntut kita mengetahui di mana $O$ duduk di dalam tabungnya.

\textbf{6.} $\partial f_1'/\partial D = 1/4 + d^2/4D^2 = 0.3125$,
$\partial f_1'/\partial d = -d/2D = -0.25$; sehingga
$u(f_1')^2 = (0.3125 \times 0.2)^2 + (0.25 \times 0.5)^2 =
\qty{0.0195}{cm^2}$ dan $u(f_1') = \qty{0.14}{cm}$:
$f_1' = \qty{15.00\pm0.14}{cm}$.

\textbf{7.} $p = (80 \mp 40)/2 = 20$ atau \qty{60}{cm}: $\gamma =
-60/20 = -3$ (slide dekat \omterm{def:b1:mirrors-thin-lenses:lens}{lensanya}, diperbesar: itulah proyektornya) dan
$\gamma = -1/3$.

\textbf{8.} $D_{\min} = 4f_1' = \qty{60.0}{cm}$, $\gamma = -1$.

\textbf{9.} $\gamma = -1.80/0.036 = -50$.

\textbf{10.} $\overline{OA'} = f_1'(1 - \gamma) = 15 \times 51 =
\qty{765}{cm} = \qty{7.65}{m}$; $\overline{OA} = \overline{OA'}/\gamma =
\qty{-15.3}{cm}$.

\textbf{11.} $\gamma < 0$: \omterm{def:b1:mirrors-thin-lenses:image}{bayangan nyatanya} diputar $180^\circ$
terhadap sumbunya, sehingga slidenya dimasukkan setelah diputar
$180^\circ$ (terbalik atas--bawah dan kiri--kanan).

\textbf{12.} $\overline{OA'} = \qty{1000}{cm}$: $1/\overline{OA} = 1/1000
- 1/15$, $\overline{OA} = \qty{-15.23}{cm}$: \omterm{def:b1:mirrors-thin-lenses:lens}{lensanya} bergerak \qty{0.7}{mm}
mendekati slidenya. $\gamma = 1000/(-15.23) = -65.7$: gambarnya
selebar \qty{2.36}{m}.

\textbf{13.} Dengan menurunkan Descartes: $-\dd\overline{OA'}/\overline{OA'}^2
+ \dd\overline{OA}/\overline{OA}^2 = 0$, sehingga $\dd\overline{OA'} =
(\overline{OA'}/\overline{OA})^2\,\dd\overline{OA} = \gamma^2\,\dd\overline{OA}$.

\textbf{14.} $2500 \times \qty{0.10}{mm} = \qty{25}{cm}$: \omterm{def:b1:mirrors-thin-lenses:image}{bayangannya}
hanyut keluar fokus seiring memanasnya film --- karena itu slide berbingkai
kaca, atau autofokus yang mengikuti filmnya.

\textbf{15.} Segitiga sebangun dari bukaannya ke titik \omterm{def:b1:mirrors-thin-lenses:image}{bayangannya}:
kekaburannya $= 30 \times 0.20/7.65 = \qty{0.78}{mm}$; dari \qty{5}{m} ini
membentang $\qty{1.6e-4}{rad} < \qty{3e-4}{rad}$: tak disadari.

\textbf{16.} Luas slidenya $8.64\times10^{-4}\ \unit{m^2}$: $E_{\text{slide}}
= 2.0/8.64\times10^{-4} = \qty{2.3e3}{W/m^2}$; luas layarnya $\gamma^2$
kali lebih besar, $\qty{2.16}{m^2}$: $E_{\text{layar}} = \qty{0.93}{W/m^2}$;
dayanya sama, luasnya bernisbah $\gamma^2 = 2500$.

\textbf{17.} $\overline{O_2A'} = \qty{1200}{cm}$: $1/\overline{O_2A} =
1/1200 + 1/30 = 0.03417$, $\overline{O_2A} = \qty{+29.3}{cm}$: \omterm{def:b1:mirrors-thin-lenses:image}{bayangan}
antaranya terletak \qty{29.3}{cm} \emph{di luar} $L_2$ --- sebuah benda
maya bagi $L_2$ (sinar dari $L_1$ memusat ke arahnya).

\textbf{18.} $\overline{O_1A_1'} = 8.0 + 29.3 = \qty{37.3}{cm}$;
$1/\overline{O_1A} = 1/37.3 - 1/15 = -0.0398$, $\overline{O_1A} =
\qty{-25.1}{cm}$.

\textbf{19.} $\gamma_1 = 37.3/(-25.1) = -1.49$; $\gamma_2 = 1200/29.3 =
41.0$; $\gamma = -60.9$.

\textbf{20.} Lebarnya $60.9 \times 36 = \qty{2.19}{m}$. $L_1$ sendirian:
$\gamma = 1 - \overline{OA'}/f_1' = 1 - 1208/15 = -79.5$, lebarnya
\qty{2.86}{m}.

\textbf{21.} Berkas bersudut $\alpha$: $L_1$ memfokuskannya pada ketinggian $h_1 = f_1'\alpha$
di bidang fokusnya, yang berjarak $\overline{O_2A_1} = f_1' - e$ dari $L_2$;
$L_2$ membayangkan titik ini di $\overline{O_2A'}$ dengan $1/\overline{O_2A'}
= 1/f_2' + 1/(f_1' - e)$ dan memperbesarnya sebesar $\gamma_2 = \overline{O_2A'}/(f_1'
- e)$, sehingga $h = \gamma_2 f_1'\alpha$ dan $f_{\mathrm{eq}} = \gamma_2 f_1'
= \dfrac{f_1' f_2'}{f_1' + f_2' - e}$, yang kebalikannya adalah rumus
yang dinyatakan. $1/f_{\mathrm{eq}} = 1/15 - 1/30 + 8/450 = \qty{0.0511}{cm^{-1}}$,
$f_{\mathrm{eq}} = \qty{19.6}{cm}$.

\textbf{22.} $2.19/2.86 = 0.77$ dan $15/19.6 = 0.77$. Hanya hampir,
karena slidenya berada pada jarak berhingga dan kedua lemparannya diukur
dari \omterm{def:b1:mirrors-thin-lenses:lens}{lensa} yang berbeda; untuk benda di tak hingga nisbahnya akan eksak.

\textbf{23.} Benda maya pada $p > 0$ di luar \omterm{def:b1:mirrors-thin-lenses:lens}{lensa divergen}:
$1/\overline{O_2A'} = 1/p + 1/f_2' = 1/p - 1/\abs{f_2'}$, positif
(\omterm{def:b1:mirrors-thin-lenses:image}{bayangan nyata}) jika dan hanya jika $p < \abs{f_2'}$.

\textbf{24.} Daya yang sama pada gambar yang luasnya $(2.19/2.86)^2 = 0.59$
kali lebih kecil: iradiansinya $\times 1.7$.

\textbf{25.} Pengubah itu mengalikan jarak fokusnya dengan
$f_{\mathrm{eq}}/f_1' = 1.3$; pada lemparan tertentu gambarnya $1.3$
kali lebih kecil pada tiap arah dan $1.7$ kali lebih cerah --- sebuah
\omterm{def:b1:mirrors-thin-lenses:lens}{lensa} lempar jauh.
\end{solution}
